\section{问题二：整个烘干过程}
\label{sec:q2}

\subsection{模型的建立}
\label{sec:q2model}

\subsubsection{控制方程与变物性闭合}

问题二把物性公式换成附录 3，并将求解时段扩展到整个烘干过程。把表~\ref{tab:property-sets} 的附录 3 公式代入式\eqref{eq:heat}、\eqref{eq:mass}，得到本问的控制方程组

\begin{equation}
\rho(\C)\cp(\C)\frac{\partial \T}{\partial \t}
=\frac{1}{\rpos}\frac{\partial}{\partial \rpos}\left(\rpos\,\kcond(\C)\frac{\partial \T}{\partial \rpos}\right),
\qquad
\frac{\partial \C}{\partial \t}
=\frac{1}{\rpos}\frac{\partial}{\partial \rpos}\left(\rpos\,\Ddiff(\C,\T)\frac{\partial \C}{\partial \rpos}\right),
\label{eq:q2system}
\end{equation}

其中

\begin{equation}
\rho=650+128\C,\qquad
\cp=1450+2736\wfrac,\qquad
\kcond=0.21+0.38\wfrac,\qquad
\Ddiff=2.4\times10^{-3}e^{-0.45/\C}e^{-3850/\T},
\qquad \wfrac=\frac{\C}{1+\C}.
\label{eq:app3}
\end{equation}

式中四个物性\textbf{都随状态变化}，且 $\Ddiff$ 同时依赖 $\C$ 与 $\T$（$\T$ 用开尔文）。与问题一的常数物性相比，本问的方程在数学性质上有三点变化：热容量 $\rho\cp$ 出现在时间导数之前且随时间变化，方程不再是线性抛物型；扩散系数含有对温度敏感的指数因子，使温度场与水分场形成 $\T\to\Ddiff$ 的耦合；而含水率又通过 $\rho,\cp,\kcond$ 反过来影响热量输运，形成 $\C\to(\rho,\cp,\kcond)$ 的反向耦合。两条耦合路径都不含``源项''，因此耦合是\textbf{弱耦合}：它不改变方程的抛物型性质，只改变系数。

定解条件沿用式\eqref{eq:initial}$\sim$\eqref{eq:robin}，物性按式\eqref{eq:app3} 取值。表面热边界仍为纯对流（假设 A5），表面质边界仍为题给的浓度差驱动形式（假设 A7）。

\subsubsection{关于两阶段的处理}

题目指出烘干过程包含预热平衡与恒温干燥两个阶段、且两阶段参数有所不同，但同时又要求``为简化问题，相关经验公式统一采用附录 3 中的公式''。也就是说，\textbf{参数集上的阶段差异已被题面显式抹平}，剩下的差别只有环境条件 $\Tinf(\t)$、$\Cinf(\t)$ 的时程：附件 1 给出前 $4$~h 的实测序列，此后环境按假设 B2 保持末值。

据此本文\textbf{不引入人为的阶段切换判据}。理由有二：一是切换时刻需要额外的判据（例如``表面温度达到环境温度的 $95\%$''），而该判据的阈值题目未给，引入它等于凭空增加一个参数；二是本问的方程形式与参数集在两阶段相同，强行分段只会引入一个不连续的接点，反而可能造成数值上的伪振荡。相应地，本文只在结果分析中以环境条件的时间历程来叙述``预热段''与``恒温段''，而不声称两者之间存在可辨识的数学分界。

\subsection{模型的求解}
\label{sec:q2analysis}

\subsubsection{非线性与刚性的处理}

式\eqref{eq:q2system} 的刚性来自两方面：一是表面加密网格的最小单元（$\Nr=2561$ 时无量纲宽度约 $2.4\times10^{-6}$），它决定了显式格式的稳定性上界约为 $10^{-8}$~s 量级；二是扩散系数在干燥过程中跨越两个数量级以上（表~\ref{tab:scale}），使方程的谱半径随时间剧烈变化。因此本文采用第~\ref{sec:time} 节的隐式 BDF 积分，物性在每次右端求值时按当前 $(\C,\T)$ 重新计算，非线性由积分器内部的 Newton 迭代处理，不下做滞后或外推。

从计算量看本问并不困难：正式计算（$\Nr=2561$）的墙钟时间约 \QTwoWallS~s，远低于预算，故本文没有采用任何降精度措施（无粗网格、无放宽容差、无输出平滑）。

\subsubsection{时间尺度分离与内部扩散控制}

本问的物理图像可以用表~\ref{tab:scale} 说明。表中的量都在末端环境温度 $323.315$~K 下用式\eqref{eq:app3} 实算：$\tau_C={\Rzero}^2/\Ddiff$ 是水分扩散的特征时间，$\hmass\Rzero/\Ddiff$ 是表面传质阻力与内部扩散阻力之比，${\Rzero}^2\rho\cp/\kcond$ 是导热特征时间。

\begin{table}[!ht]
  \centering
  \caption{附录 3 物性在末端环境温度 $323.315$~K 下随含水率的变化（均为由题给公式实算的量级估计）}
  \label{tab:scale}
  \scriptsize
  \begin{tabular}{ccccccc}
    \toprule
    $\C$ / (kg/kg) & $\wfrac$ & $\rho$ / (kg/m$^3$) & $\cp$ / (J/(kg$\cdot$K))
      & $\kcond$ / (W/(m$\cdot$K)) & $\Ddiff$ / (m$^2$/s)
      & $\tau_C$ / h \\
    \midrule
    $2.55$（初值） & $0.718$ & $976.4$ & $3415.3$ & $0.4830$ & $1.36\times10^{-8}$ & $8.20$ \\
    $1.50$ & $0.600$ & $842.0$ & $3091.6$ & $0.4380$ & $1.20\times10^{-8}$ & $9.28$ \\
    $1.00$ & $0.500$ & $778.0$ & $2818.0$ & $0.4000$ & $1.03\times10^{-8}$ & $10.78$ \\
    $0.50$ & $0.333$ & $714.0$ & $2362.0$ & $0.3367$ & $6.57\times10^{-9}$ & $16.90$ \\
    $0.15$（判据） & $0.130$ & $669.2$ & $1806.9$ & $0.2596$ & $8.05\times10^{-10}$ & $138.03$ \\
    $0.05$（末期表面） & $0.048$ & $656.4$ & $1580.3$ & $0.2281$ & $2.00\times10^{-12}$ & $5.6\times10^{4}$ \\
    \bottomrule
  \end{tabular}
\end{table}

由表~\ref{tab:scale} 可得三条结论。

\textbf{其一，热量比水分快两个数量级。} 导热特征时间在初值处仅 $0.77$~h，远小于干燥时长；实测也印证了这一点——$t=2$~h 时中心与表面温度已分别达到 $48.45~^\circ$C 与 $49.00~^\circ$C，$t=6$~h 时全场均为 $50.1645~^\circ$C，此后直到烘干结束温度场始终保持在 $50.165~^\circ$C 的环境温度上（表~\ref{tab:q2-temperature} 与图~\ref{fig:q1T}）。因此\textbf{干燥过程的时间尺度完全由传质决定}，这一判断也是后文把温度场近似为``准均匀''来解释结果的依据。

\textbf{其二，干燥阻力由内部扩散控制，且随含水率下降急剧增大。} 比值 $\hmass\Rzero/\Ddiff$ 由初值处的 $1.18$（此时表面阻力与内部阻力同量级）增大到判据值处的 $19.9$、末期表面处的近 $10^4$。这表示后期\textbf{表面水分被气流迅速带走、内部供给跟不上}：表~\ref{tab:q2-moisture} 显示 $24$~h 时表面含水率已降到 $0.0663$~kg/kg（接近环境参考值 $0.04986$），而轴心仍有 $0.2371$~kg/kg；此后 $24$~h 内轴心缓慢降到 $0.1500$，表面只在 $0.0525\sim0.0536$ 附近微降。

\textbf{其三，干燥后期在表面附近形成陡峭的浓度边界层。} 扩散特征时间从 $8.2$~h 增大到 $138$~h 乃至更大，意味着水分``再分布''的速度远慢于表面失水的速度，故含水率的径向梯度在末期集中在靠近表面的薄层内。这正是本文采用表面二次加密网格（式\eqref{eq:gradedmesh}）的定量理由：若用均匀网格，解析这一薄层所需的节点数将高出一到两个数量级。

\subsubsection{算法的分段与实现}

求解流程为：先按式\eqref{eq:gradedmesh} 生成网格与体积权重，初值取式\eqref{eq:initial}，然后在附件 1 的每个测量时刻（$0,60,\ldots,14400$~s）以及外推段内分段积分，每段重启积分器。这样做的目的是避免单个自适应步长跨越由式\eqref{eq:interp} 引入的折点，同时让外推接点（$t=14400$~s）处参数的切换发生在段的边界上。终点判据与问题三一致（全场最大含水率低于 $0.15$~kg/kg，见式~\eqref{eq:maxcrit}），使本问的求解区间与问题三完全对应，便于两问互相校验：实测两者的事件时刻在 $1$~s 的输出网格上完全一致。

\subsection{结果与分析}

按题目表 3、表 4 的格式整理的结果见表~\ref{tab:q2-temperature}、表~\ref{tab:q2-moisture}。

\geninput{tab/q2_temp}{表 3 由 \texttt{make tables} 从 \texttt{results/tables/q2\\_temp.csv} 生成。}
\geninput{tab/q2_moist}{表 4 由 \texttt{make tables} 从 \texttt{results/tables/q2\\_moist.csv} 生成。}

\textbf{温度：全场在前 $2$~h 内完成升温，随后保持与环境同温。} 表~\ref{tab:q2-temperature} 给出 $0.5\sim3$~h 的六个时刻：$0.5$~h 时中心与表面分别为 $32.19~^\circ$C 与 $35.41~^\circ$C（径向温差 $3.22~^\circ$C，与问题一 $1800$~s 的结果一致），$1$~h 时升到 $40.38~^\circ$C 与 $43.00~^\circ$C，$2$~h 时已达 $48.45~^\circ$C 与 $49.00~^\circ$C，$3$~h 时全场都在 $50~^\circ$C 附近且径向差不足 $0.12~^\circ$C。此后温度场与环境的差别小于 $10^{-3}~^\circ$C，故整个干燥过程可以视为\textbf{等温过程}，这与第~\ref{sec:q2analysis} 节的尺度分析一致。


\textbf{水分：呈现出``表层快速下降—内部缓慢跟进—末期整体收敛''的三段式规律。} 由表~\ref{tab:q2-moisture} 与图~\ref{fig:q1C}：

\begin{itemize}
  \item \textbf{第 1 段（$0\sim3$~h）}：失水集中在表层。$3$~h 时表面已降到 $1.0081$~kg/kg（失去 $60\%$ 的初始含水率），而轴心仍有 $1.7662$~kg/kg，$0.5$~cm 处为 $1.7165$~kg/kg、$1.5$~cm 处为 $1.3333$~kg/kg。剖面的形状是``内部平台 $+$ 近表面陡降''。
  \item \textbf{第 2 段（$3\sim24$~h）}：干燥前沿向内部推进，内部含水率开始快速下降。$6$~h 时轴心为 $1.0156$~kg/kg，$12$~h 降到 $0.4548$，$24$~h 为 $0.2371$；同期表面已逼近环境参考值（$24$~h 为 $0.0663$~kg/kg）。这一段的表观干燥速率最大（$12\sim24$~h 内轴心每小时下降约 $0.018$~kg/kg）。
  \item \textbf{第 3 段（$24$~h 之后）}：全场的失水速率显著变缓，含水率剖面整体接近水平并缓慢逼近判据值。$48$~h 时中心为 $0.1615$、表面为 $0.0536$~kg/kg，中心每小时仅下降约 $0.003$~kg/kg，而表面上 $24$~h 内只下降 $0.0127$~kg/kg。
\end{itemize}


\textbf{三段式规律的物理原因正是表~\ref{tab:scale} 所示的扩散系数变化。} 第 1 段中 $\Ddiff$ 仍较大（$10^{-8}$~m$^2$/s 量级），表层水分可以通过扩散迅速补充并蒸发，故表现为``表面快速失水、内部基本不动''；第 2 段中含水率整体下降到 $0.5$ 量级，$\Ddiff$ 降到 $6.57\times10^{-9}$~m$^2$/s，内部补充变慢，于是整个断面的含水率开始同步下降；第 3 段中含水率接近 $0.15$，$\Ddiff$ 已降到 $8\times10^{-10}$~m$^2$/s，扩散特征时间达到 $138$~h，远大于当时的剩余干燥时间，因此\textbf{含水率只能以极慢的速度``挤''向判据值}。这解释了为什么最后 $10\%$ 的含水率下降需要消耗整个干燥过程 $60\%$ 的时间。

全过程烘干时长为 \textbf{\QTwoEventH~h}（\QTwoEventS~s），此时全场最大含水率为 $0.15$~kg/kg（判据值），交付文件写到第 \QTwoLastRowS~s。需要说明的是，本问的终点时刻与问题三完全相同，这不是巧合，而是因为两问的方程、参数集、几何与判据完全一致，只有输出步长不同（$1$~s 与 $60$~s）。因此本问的终点时刻同时可视为对问题三结果的交叉印证。

最后，关于
